% from ega2-1-fr.tex
\oldpage[II]{99}

En effet, on a la suite exacte
\[
  0\to\mathrm{H}^1(Y,f_*(\mathcal{F}))\to
  \mathrm{H}^1(X,\mathcal{F})\to
  \mathrm{H}^0(Y,\mathrm{R}^1f_*(\mathcal{F}))
\]
des termes de bas degré de la suite spectrale de Leray, et la conclusion
résulte de (\hyperref[II.5.2.2-fr]{5.2.2}).

\begin{remark}[5.2.4]
\label{II.5.2.4-fr}
Au chapitre III, § 1, nous prouverons que si $X$ est affine, on a
$\mathrm{H}^i(X,\mathcal{F})=0$ pour tout $i>0$ et tout
$\mathcal{O}_X$-Module quasi-cohérent $\mathcal{F}$.
\end{remark}

\subsection{Morphismes quasi-projectifs.}
\label{subsection:II.5.3-fr}

\begin{definition}[5.3.1]
\label{II.5.3.1-fr}
On dit qu'un morphisme $f:X\to Y$ est \emph{quasi-projectif}, ou que $X$
(considéré comme $Y$-préschéma au moyen de $f$) est \emph{quasi-projectif
sur $Y$}, ou que $X$ est un \emph{$Y$-schéma quasi-projectif}, si $f$ est de
type fini et s'il existe un $\mathcal{O}_X$-Module inversible qui est
$f$-ample.
\end{definition}

On notera que cette notion n'est pas locale sur $Y$~: les contre-exemples de
Nagata [26] et de Hironaka montrent que, même si $X$ et $Y$ sont des schémas
algébriques non singuliers sur un corps algébriquement clos, tout point de
$Y$ peut avoir un voisinage affine $U$ tel que $f^{-1}(U)$ soit
quasi-projectif sur $U$, sans que $f$ soit quasi-projectif.

On notera qu'un morphisme quasi-projectif est nécessairement séparé
(\hyperref[II.4.6.1-fr]{4.6.1}). Lorsque $Y$ est quasi-compact, il revient au
même de dire que $f$ est quasi-projectif ou que $f$ est de type fini et qu'il
existe un $\mathcal{O}_X$-Module très ample relativement à $f$
(\hyperref[II.4.6.2-fr]{4.6.2} et \hyperref[II.4.6.11-fr]{4.6.11}). De plus~:

\begin{proposition}[5.3.2]
\label{II.5.3.2-fr}
Soit $Y$ un schéma quasi-compact ou un préschéma dont l'espace sous-jacent
est noethérien, et soit $X$ un $Y$-préschéma. Les conditions suivantes sont
équivalentes~:
\begin{enumerate}
  \item[\rm{a)}] $X$ est un $Y$-schéma quasi-projectif.
  \item[\rm{b)}] $X$ est de type fini sur $Y$, et il existe un
  $\mathcal{O}_Y$-Module quasi-cohérent de type fini $\mathcal{E}$ tel que
  $X$ soit $Y$-isomorphe à un sous-préschéma de $\mathbf{P}(\mathcal{E})$.
  \item[\rm{c)}] $X$ est de type fini sur $Y$, et il existe une
  $\mathcal{O}_Y$-Algèbre graduée quasi-cohérente $\mathcal{S}$ telle que
  $\mathcal{S}_1$ soit de type fini et engendre $\mathcal{S}$ et que $X$
  soit $Y$-isomorphe à un sous-préschéma induit sur un ouvert partout dense
  de $\operatorname{Proj}(\mathcal{S})$.
\end{enumerate}
\end{proposition}

Cela résulte aussitôt de la remarque précédente et de
(\hyperref[II.4.4.3-fr]{4.4.3}), (\hyperref[II.4.4.6-fr]{4.4.6}) et
(\hyperref[II.4.4.7-fr]{4.4.7}).

On notera que lorsque $Y$ est un préschéma noethérien, on peut, dans les
conditions b) et c) de (\hyperref[II.5.3.2-fr]{5.3.2}), supprimer
l'hypothèse que $X$ est de type fini sur $Y$, qui est alors automatiquement
vérifiée (\hyperref[I.6.3.5-fr]{I, 6.3.5}).

\begin{corollary}[5.3.3]
\label{II.5.3.3-fr}
Soit $Y$ un schéma quasi-compact tel qu'il existe un $\mathcal{O}_Y$-Module
ample $\mathcal{L}$ (\hyperref[II.4.5.3-fr]{4.5.3}). Pour qu'un $Y$-schéma
$X$ soit quasi-projectif, il faut et il suffit qu'il soit de type fini sur
$Y$ et isomorphe à un sous-$Y$-schéma d'un fibré projectif de la forme
$\mathbf{P}_Y^r$.
\end{corollary}

En effet, si $\mathcal{E}$ est un $\mathcal{O}_Y$-Module quasi-cohérent de
type fini, $\mathcal{E}$ est isomorphe à un quotient d'un
$\mathcal{O}_Y$-Module
$\mathcal{L}^{\otimes(-n)}\otimes_{\mathcal{O}_Y}\mathcal{O}_Y^k$
(\hyperref[II.4.5.5-fr]{4.5.5}), donc $\mathbf{P}(\mathcal{E})$ est
isomorphe à un sous-schéma fermé de $\mathbf{P}_Y^{k-1}$
(\hyperref[II.4.1.2-fr]{4.1.2} et \hyperref[II.4.1.4-fr]{4.1.4}).

\begin{proposition}[5.3.4]
\label{II.5.3.4-fr}
\begin{enumerate}
  \item[\rm{(i)}] Un morphisme quasi-affine de type fini (et en particulier
  une immersion quasi-compacte, ou un morphisme affine de type fini) est
  quasi-projectif.
  \item[\rm{(ii)}] Si $f:X\to Y$ et $g:Y\to Z$ sont quasi-projectifs et si
  $Z$ est quasi-compact, $g\circ f$ est quasi-projectif.

